\subsection{根系构造的预备知识}

设 V 是一个维数 \(l \geqslant 1\) 的实向量空间，并配备标量积 \((x|y)\)。设 L 是 V 的离散子群，A 是一个有限的数集，且其中每个数为 \textgreater0；R 是由这样的元素组成的集合：\(\alpha \in L\)，且 \((\alpha|\alpha) \in \Lambda\)。假设 R 生成 V，并且对于 R 中任意一对点 \((\alpha,\beta)\)，数 \(2\frac{(\alpha|\beta)}{(\alpha|\alpha)}\) 是整数。则 R 是 V 中的根系。事实上，R 显然满足 \((RS_{I})\)。令 \(\alpha \in R\)：令 \(s_{\alpha}\) 为正交反射 \(x \mapsto x - 2\frac{(x|\alpha)}{(\alpha|\alpha)}\alpha\)；于是若 \(\beta \in R\)，有 \(2\frac{(\beta|\alpha)}{(\alpha|\alpha)} \in Z\)，故 \(s_{\alpha}(\beta) \in L\)，且 \(\|s_{\alpha}(\beta)\| = \| \beta \|\)，所以 \(s_{\alpha}(\beta) \in R\)；因此 R 满足 \((RS_{II})\) 和 \((RS_{III})\)，并且若 A 不含形如 \(\lambda\) 和 \(4\lambda\) 的两个数，则 R 是约化的。

现在令 V 为 \(E = R^{n}\) 的一个子空间。令 \((\varepsilon_{1}, \ldots, \varepsilon_{n})\) 为 E 的典范基；在 E 上取使此基成为标准正交基的标量积，并借助此标量积将 \(E^{*}\) 与 E（分别将 \(V^{*}\) 与 V）认同。定义 E 的子群 \(L_{0}, L_{1}, L_{2}, L_{3}\) 如下：

\begin{enumerate}
\def\labelenumi{\arabic{enumi})}
\item
  \(L_{0}\) 是以 \((\varepsilon_{i})\) 为基的 Z-模。我们有 \((\alpha|\beta)\in\mathbf{Z}\)，其中 \(\alpha,\beta\in L_{0}\)。满足条件 \(\alpha\in L_{0}\) 且 \((\alpha|\alpha)=1\) 的向量是 \(\pm\varepsilon_{i}\;(1\leqslant i\leqslant n)\)；满足 \((\alpha|\alpha)=2\) 的向量是 \(\pm\varepsilon_{i}\pm\varepsilon_{j}\)，其中 i\textless j（\(\pm\) 中的两个符号，在 \(\pm\varepsilon_{i}\pm\varepsilon_{j}\) 中彼此独立选取；本节其余部分均采用类似约定）。
\item
  \(L_{1}\) 是 \(L_{0}\) 的 Z-子模，由 \(x = \sum_{i=1}^{n} \xi_{i} \varepsilon_{i} \in L_{0}\) 组成，其中 \(\sum_{i=1}^{n} \xi_{i}\) 为偶数；由于 \(\xi_{i}\) 与 \(\xi_{i}^{2}\) 具有相同的奇偶性，这等价于要求 \((x|x)\) 为偶数。令 \(L_{1}^{\prime}\) 为 \(L_{1}\) 的子模，由 \(\varepsilon_{i} \pm \varepsilon_{j}\) 生成；我们有 \(\sum_{i=1}^{n} \xi_{i} \varepsilon_{i} \equiv (\sum_{i=1}^{n} \xi_{i}) \varepsilon_{n} \pmod{L_{1}^{\prime}}\)，且由于 \(2\varepsilon_{n} = (\varepsilon_{1} + \varepsilon_{n}) - (\varepsilon_{1} - \varepsilon_{n}) \in L_{1}^{\prime}\)，可知 \(L_{1}^{\prime} = L_{1}\)。由于 \(L_{0}\) 由 \(L_{1}\) 和 \(\varepsilon_{1}, L_{0}/L_{1}\) 生成，后者同构于 Z/2Z。
\item
  \(\mathrm{L}_2 = \mathrm{L}_0 + \mathbf{Z}\left(\frac{1}{2} \sum_{i=1}^{n} \varepsilon_i\right)\)。显然，元素 \(x = \sum_{i=1}^{n} \xi_i \varepsilon_i\) 属于 \(\mathrm{V}\)，它属于 \(\mathrm{L}_2\) 当且仅当
\end{enumerate}

\[
2 \xi_ {i} \in \mathbf {Z}, \quad \xi_ {i} - \xi_ {j} \in \mathbf {Z} \text {   for   all   } i \text {   and   } j.\tag{6}
\]

由于 \((\varepsilon_k|_{\frac{1}{2}}\sum_{i=1}^{n}\varepsilon_i) = \frac{1}{2}\) 对所有 \(k\) 成立，且 \(\| \frac{1}{2}\sum_{i=1}^{n}\varepsilon_i\|^2 = \frac{n}{4}\)，有 \((\alpha|\beta) \in \frac{1}{2}\mathbf{Z}\) 对于 \(\alpha, \beta \in L_2\) 成立，其中 \(n\) 为偶数。群 \(L_2 / L_0\) 同构于 \(\mathbf{Z}/2\mathbf{Z}\)。

\begin{enumerate}
\def\labelenumi{\arabic{enumi})}
\setcounter{enumi}{3}
\tightlist
\item
  \(L_{3} = L_{1} + \mathbf{Z}\left(\frac{1}{2} \sum_{i=1}^{n} \varepsilon_{i}\right)\)。若 n 是 4 的倍数，则 \(L_{3}\) 是由 \(\sum_{i=1}^{n} \xi_{i} \varepsilon_{i}\) 组成的集合，这些元素满足 (6) 以及条件 \(\sum_{i=1}^{n} \xi_{i} \in 2\mathbf{Z}\)；在此情形，\((\alpha | \beta) \in \mathbf{Z}\) 对所有 \(\alpha, \beta \in L_{3}\) 成立。
\end{enumerate}

立即可见，与 \(L_{0}\)（分别为 \(L_{1}, L_{2}\)）关联的 E 的子群是 \(L_{0}\)（分别为 \(L_{2}, L_{1}\)）。与 \(L_{3}\) 关联的 E 的子群是

\[
x = \sum_ {i = 1} ^ {n} \xi_ {i} \varepsilon_ {i} \in L _ {2}
\]

满足 \((x|\frac{1}{2}\sum_{i = 1}^{n}\varepsilon_i)\in \mathbf{Z}\)，即满足 \(\sum_{i = 1}^{n}\xi_{i}\in 2\mathbf{Z}\)；因此若 \(n\equiv 0\)（模 4），这个关联子群就是 \(\mathrm{L}_3\)。

阿贝尔群 \(\mathrm{L}_2 / \mathrm{L}_1\) 的阶为 4，因而同构于 \(\mathbf{Z} / 4\mathbf{Z}\) 或 \(\mathbf{Z} / 2\mathbf{Z} \times \mathbf{Z} / 2\mathbf{Z}\)（《代数》第七章第 4 节第 6 号定理 3）。若 \(n\) 为奇数，

\[
p (\frac {1}{2} \sum_ {i = 1} ^ {n} \varepsilon_ {i}) \in \mathrm{L} _ {1} \Leftrightarrow p \equiv 0 (\mathrm{mod.} 4)
\]

所以 \(\mathrm{L}_2 / \mathrm{L}_1\) 是 4 阶循环群。若 \(n\) 为偶数，

\[
p (\frac {1}{2} \sum_ {i = 1} ^ {n} \varepsilon_ {i}) \in \mathrm{L} _ {1} \Leftrightarrow p \equiv 0 (\mathrm{mod.} 2)
\]

所以 \(L_{2}/L_{1}\) 含有两个不同的 2 阶元素，并同构于 \(Z/2Z \times Z/2Z\)。

在下面九个号以及各图版中，我们将一直使用这些记号。对于定理 3 中的每一种 Dynkin 图，我们将明确描述：

\begin{enumerate}
\def\labelenumi{(\Roman{enumi})}
\item
  一个根系 R 及其根的数目。
\item
  R 的一组基 B 及相应的正根。基 B 将用整数 \(1,\ldots,l\) 编号。
\item
  Coxeter 数 \(h\)（第 1 节，第 11 号）。
\item
  最高根 \(\tilde{\alpha}\)（相对于 B 所定义的序）及完备 Dynkin 图（第 3 号）。我们将在每个顶点旁标出 B 中对应的根。
\item
  对偶根系 \(\mathbb{R}^{\prime}\)、典范双线性形式以及常数 \(\gamma(\mathbb{R})\)（第 1 节，第 12 号）。
\item
  相对于 B 的基本权（第 1 节，第 10 号）。
\item
  正根之和。
\item
  群 \(P(R), Q(R), P(R)/Q(R)\) 以及连接指标（第 1 节，第 9 号）。
\item
  \(\mathrm{W}(\mathrm{R})\) 的指数（第五章第 6 节第 2 号定义 2）。在 \(\mathbf{A}_l, \mathbf{B}_l, \mathbf{C}_l\) 和 \(\mathbf{D}_l\) 的情形，我们确定其对称不变量。
\item
  W(R) 的阶（以及在某些情形下它的结构）。
\item
  群 \(A(R)/W(R)\)、它在 Dynkin 图上的作用，以及将 B 变为 -B 的元素 \(w_{0}\)（属于 \(W(R)\)）。
\item
  \(\mathrm{P}(\mathrm{R}^{\prime})/\mathrm{Q}(\mathrm{R}^{\prime})\) 在完备 Dynkin 图上的作用，以及 \(\mathrm{A}(\mathrm{R})/\mathrm{W}(\mathrm{R})\) 在 \(\mathrm{P}(\mathrm{R}^{\prime})/\mathrm{Q}(\mathrm{R}^{\prime})\) 上的作用。
\end{enumerate}

对于定理 3 中的每个 Dynkin 图，这些数据将汇集于图版 I 至 IX 中，并按上述统一方式排列。我们还给出：

\begin{enumerate}
\def\labelenumi{(\Roman{enumi})}
\setcounter{enumi}{12}
\tightlist
\item
  Cartan 矩阵，由它按第 2 号所述方式导出 Dynkin 图。
\end{enumerate}

